\section{Obstacle envelopes}
\label{sec:obstacle-envelopes}

An obstacle selects the largest potential lying below it. The local upper-envelope theorem gives the construction once the competitors have a common upper bound; normalized compactness supplies that bound for obstacles finite almost everywhere.

A smooth real $(1,1)$-form $\theta$ is called \emph{big} if there are $\rho\in\PSH(X,\theta)$ and $\delta>0$ such that
\[
 \theta+\ddc\rho\geq\delta\omega.
\]
Here $\omega$ is any Hermitian metric. Changing that metric changes $\delta$, not the condition. Neither this definition nor the existence of $\theta$-plurisubharmonic functions requires $\theta$ to be closed.

For any smooth real $(1,1)$-form $\theta$ and an extended-real obstacle $f$ on $X$, defined at least outside a pluripolar set, write
\[
 P_\theta(f)=\left(\sup\{v\in\PSH(X,\theta):v\leq f
                      \ \hbox{q.e.}\}\right)^*.
\]
For $u,v\in\PSH(X,\theta)$, set $P_\theta(u,v)=P_\theta(\min(u,v))$. The star denotes upper semicontinuous regularization, and the empty supremum is $-\infty$. Quasi-everywhere means outside a pluripolar set. This convention also gives meaning to differences of quasi-plurisubharmonic functions at common poles. For an upper semicontinuous obstacle, the pointwise definition gives the same competitors: a quasi-plurisubharmonic function lying below it almost everywhere lies below it everywhere, by the local submean inequality and upper semicontinuity of the obstacle.

The local envelope theorem from \cite{SGPT} states that the supremum of a nonempty, locally uniformly upper-bounded family of plurisubharmonic functions has a plurisubharmonic upper semicontinuous regularization. This regularization differs from the supremum only on a pluripolar set. To retain the lower bound $-\theta$, fix a constant positive form $\kappa$ in local coordinates. For each $\epsilon>0$, shrink the neighborhood and choose a constant form $\theta'$ with $-\epsilon\kappa\leq\theta-\theta'\leq\epsilon\kappa$. A smooth potential $r$ for $\theta'+\epsilon\kappa$ makes every competitor plus $r$ plurisubharmonic. The local envelope theorem then gives $\theta+\ddc P_\theta(f)\geq-2\epsilon\kappa$ there. Letting $\epsilon$ tend to zero proves the exact bound. The same local theorem shows that regularization changes the supremum only on a pluripolar set.

Here is a useful way to check the upper bound for an unbounded obstacle. If $f\leq A$ for a real constant $A$ on a Borel set $E$ of positive $\mathrm{d}V$-measure and $h\leq f$ is a competitor, normalize $\widetilde h=h-\sup_Xh$. By \cref{thm:normalized-compactness},
\[
 (\sup_Xh-A)\operatorname{vol}_{\mathrm{d}V}(E)
 \leq\int_E|\widetilde h|\,\mathrm{d}V\leq C_\theta
 \quad\hbox{whenever }\sup_Xh>A.
\]
Thus the competitors have a common upper bound. In particular this applies to any measurable obstacle finite almost everywhere, by taking one of its finite sublevel sets of positive volume.

If $\PSH(X,\theta)$ is nonempty, its normalized envelope
\[
 V_\theta=P_\theta(0)
\]
is $\theta$-plurisubharmonic. Indeed, a normalized $\theta$-plurisubharmonic function is a competitor, and the family is bounded above by zero. Moreover,
\[
 u-\sup_Xu\leq V_\theta\qquad(u\in\PSH(X,\theta)).
\]
Thus $V_\theta$ has minimal singularities. If $\theta\geq0$, then $V_\theta=0$.

\begin{lemma}\label{lem:bounded-below-envelope}
Let $\eta$ be a smooth semipositive real $(1,1)$-form on $X$. If an obstacle $f$ is finite almost everywhere and bounded below quasi-everywhere by a constant, then $P_\eta(f)$ is a bounded $\eta$-plurisubharmonic function.
\end{lemma}
\begin{proof}
The lower constant is a competitor. The preceding finite-sublevel argument gives a common upper bound for all competitors. The local envelope theorem therefore gives a bounded envelope.
\end{proof}

\begin{proposition}
\label{prop:continuous-obstacle-envelope}
If $f$ is continuous and $\omega$ is Hermitian, then $P_\omega(f)$ is continuous. Moreover,
\[
 \|P_\omega(f)-P_\omega(g)\|_\infty\leq\|f-g\|_\infty
 \quad(f,g\in C^0(X)).
\]
\end{proposition}
\begin{proof}
The constant $\min_Xf$ is a competitor, so $p=P_\omega(f)$ is bounded. The local envelope theorem gives $p\leq f$ quasi-everywhere, and continuity of $f$ makes this inequality pointwise. Choose smooth $v_j\downarrow p$ by \cref{thm:bounded-exact-regularization}. For every $\epsilon>0$, the continuous functions $v_j-f$ decrease to a function at most zero. Compactness shows that $v_j\leq f+\epsilon$ everywhere for all large $j$: otherwise their decreasing compact $\epsilon$-superlevel sets would have nonempty intersection. Hence $v_j-\epsilon$ is a competitor and
\[
 p\leq v_j\leq p+\epsilon.
\]
Thus $v_j\to p$ uniformly, proving continuity. The last inequality follows directly from the competitors and
$P_\omega(g+c)=P_\omega(g)+c$ for constants $c$.
\end{proof}

\begin{definition}\label{def:singularity-envelope}
For $u\in\PSH(X,\theta)$, set
\[
 P_\theta[u]
 =\left(\lim_{C\to+\infty}P_\theta(u+C,0)\right)^*
 =\left(\sup\{h\in\PSH(X,\theta):h\preceq u,\ h\leq0\}\right)^*.
\]
We call $u$ an \emph{envelope fixed point} if $P_\theta[u]=u$.
\end{definition}

The two formulas agree by taking the union of the competitor families. They also agree with the definition using $V_\theta$ in place of zero: every $\theta$-plurisubharmonic competitor below zero is below $V_\theta$. The envelope is a genuine potential, since the competitors contain $u-\sup_Xu$ and are bounded above by zero. In particular, $u\preceq P_\theta[u]$. The last relation is not an assertion that the two functions differ by a bounded amount.

Directly from the competitors, $P_\theta[u]$ is unchanged when $u$ is replaced by a potential of the same ordinary singularity type, and
\[
 u\preceq v\quad\Longrightarrow\quad P_\theta[u]\leq P_\theta[v].
\]
These statements concern ordinary singularities. Replacing $\preceq$ by the multiplier-ideal order would define a different operator, which will require a separate comparison.
